\documentclass[11pt]{article}
\pdfmapfile{=lm.map}\pdfmapfile{+cm.map}\pdfmapfile{+symbols.map}
\usepackage[T1]{fontenc}
\usepackage{lmodern,amsmath,amssymb,amsthm,mathtools,booktabs,longtable,array}
\usepackage[margin=0.9in]{geometry}
\usepackage{xcolor,microtype,fancyhdr,enumitem}
\usepackage[hidelinks]{hyperref}
\definecolor{ink}{HTML}{18334B}
\pagestyle{fancy}\fancyhf{}\fancyhead[L]{\small Complex Marcus--de Oliveira}\fancyhead[R]{\small Research results, 17 October 2026}\fancyfoot[C]{\thepage}
\setlength{\headheight}{14pt}
\setlength{\parindent}{0pt}\setlength{\parskip}{6pt}
\newtheorem{theorem}{Theorem}[section]\newtheorem{lemma}[theorem]{Lemma}\newtheorem{corollary}[theorem]{Corollary}\newtheorem{proposition}[theorem]{Proposition}
\newcommand{\C}{\mathbb C}\newcommand{\R}{\mathbb R}\newcommand{\conv}{\operatorname{conv}}\newcommand{\diag}{\operatorname{diag}}\newcommand{\tr}{\operatorname{tr}}
\newcommand{\supp}{\operatorname{supp}}\newcommand{\id}{\mathrm{id}}\newcommand{\per}{\operatorname{per}}
\begin{document}
\begin{center}
{\LARGE\bfseries\color{ink} Complex Marcus--de Oliveira}\\[6pt]
{\Large Complete restricted proofs and exact certificates}\\[8pt]
17 October 2026
\end{center}

\textbf{Outcome.} This investigation does not prove or refute the unrestricted conjecture. It supplies complete proofs of two all-dimensional restricted results, a full order-four complex spectral family, exact support certificates, and a quantitative improvement to the supplied volume. The positive results are statements about actual determinants, not unproved reductions. No publication-priority claim is made.

The target remains, for every $n$, $a,b\in\C^n$ and $U\in U(n)$,
\begin{equation}\label{moc}
d_{a,b}(U):=\det(D_a+UD_bU^*)\in
C(a,b):=\conv\left\{z_\pi(a,b):=\prod_{i=1}^n(a_i+b_{\pi(i)}):\pi\in S_n\right\}.
\end{equation}
Here $D_a=\diag(a_1,\ldots,a_n)$. All convex combinations have real, nonnegative weights.

\begin{center}\begin{tabular}{p{0.27\textwidth}p{0.61\textwidth}}
\toprule
Result & Exact scope \\\midrule
Three-value theorem & Arbitrary dimension and arbitrary $a$; $b$ has values with multiplicities $(r,s,1)$. One probability law works for all numerical spectral values.\\
Small-spectrum theorem & Arbitrary dimension; for a specified generic pair $(a,\beta)$, every $U$ and all sufficiently small complex $t$ satisfy MOC for $(a,t\beta)$.\\
Explicit complex family & Full MOC for all $U(4)$ and all real $|u|\le1/20$ for the family in Section~\ref{fullfamily}; both spectra are simple when $u\ne0$.\\
Order-four cone & An explicit 20-generator certificate for nonnegative common-minor costs invariant under one transposition; sharp isoclinic bound $K^2\le4TD$.\\
Quantitative bound & In order four, the signed-coefficient norm is at most $3/2$, improving the volume's constants $24$ and $\sqrt{24}$.\\\bottomrule
\end{tabular}\end{center}


\newpage
\section{A complete three-value theorem in every dimension}

\begin{theorem}\label{three}
Let $r,s\ge0$ and $n=r+s+1$. Fix a unitary $U$ and a partition of its column labels into sets $L,H,\{c\}$ of sizes $r,s,1$. Put
\[
b_j=b_0\ (j\in L),\qquad b_j=b_1\ (j\in H),\qquad b_c=b_2.
\]
There exist nonnegative weights $\lambda_\pi$, summing to one, depending only on $U$ and this partition, such that simultaneously for every $a\in\C^n$ and every $b_0,b_1,b_2\in\C$,
\[
\det(D_a+UD_bU^*)=\sum_{\pi\in S_n}\lambda_\pi\prod_i(a_i+b_{\pi(i)}).
\]
In particular MOC holds when either spectrum has at most three distinct values and one of the three classes has size one.
\end{theorem}

The conclusion includes, for example, multiplicities $(2,2,1)$ in dimension five and $(3,3,1)$ in dimension seven. It does not require the three values to be collinear or real. The proof constructs the weights before choosing any spectral values.

\subsection{Adjacent projection laws admit one common coupling}
Write $V=U[:,H]$, $v=U[:,c]$, $P=VV^*$ and $Q=P+vv^*$. Thus $P,Q$ are orthogonal projections of ranks $s,s+1$. Define
\[
p_X=|\det U[X,H]|^2\quad(|X|=s),\qquad
q_Y=|\det U[Y,H\cup\{c\}]|^2\quad(|Y|=s+1).
\]
Cauchy--Binet gives $\sum_Xp_X=\sum_Yq_Y=1$.

\begin{lemma}\label{couple}
There exist $\mu_{X,Y}\ge0$, supported on $X\subset Y$, such that
\[
\sum_{Y\supset X}\mu_{X,Y}=p_X,\qquad
\sum_{X\subset Y}\mu_{X,Y}=q_Y.
\]
\end{lemma}
\begin{proof}
Let $x$ be the exterior product of the columns of $V$, and let $y=x\wedge v$. Their norms are one. For a family $\mathcal F$ of coordinate $s$-subsets, let $P_{\mathcal F}$ be coordinate projection, put $z=P_{\mathcal F}x$, and let $N(\mathcal F)$ be the family of $(s+1)$-subsets containing an element of $\mathcal F$.

Exterior multiplication $W_v:z\mapsto z\wedge v$ has norm at most one. Since $v$ is perpendicular to every factor of $x$, $W_v^*W_vx=x$. Also $z\wedge v$ is supported on $N(\mathcal F)$. Therefore
\[
\|P_{\mathcal F}x\|^2
=|\langle z\wedge v,y\rangle|
\le \|z\wedge v\|\,\|P_{N(\mathcal F)}y\|
\le \|P_{\mathcal F}x\|\,\|P_{N(\mathcal F)}y\|.
\]
Dividing when the first norm is nonzero and squaring proves all weighted Hall inequalities
\[
\sum_{X\in\mathcal F}p_X\le\sum_{Y\in N(\mathcal F)}q_Y.
\]
Use a finite network with capacities $p_X$ from the source to $X$, capacity $2$ on every inclusion edge $X\subset Y$, and capacity $q_Y$ from $Y$ to the sink. Every cut of capacity less than one must contain all neighbors of its source-side $X$-vertices, and the displayed inequalities exclude such a cut. The max-flow/min-cut theorem supplies a unit flow. Its inclusion-edge flows are the required $\mu_{X,Y}$. Zero marginal probabilities cause no difficulty.
\end{proof}

This is the finite adjacent-rank case of stochastic domination for projection determinantal measures; compare Lyons \cite{lyons}. The preceding argument supplies the exact finite inequalities needed here.

\subsection{Rank-one affinity turns the coupling into a convex proof}
For a rank-$s$ projection $P=VV^*$ and any $\delta\in\C$, the determinant lemma and Cauchy--Binet give
\begin{equation}\label{projexp}
\det(D_a+b_0I+\delta P)
=\sum_{|X|=s}p_X\prod_{i\in X}(a_i+b_0+\delta)\prod_{i\notin X}(a_i+b_0).
\end{equation}
For clarity, when $D=D_a+b_0I$ is invertible the left side is
\[
\det D\det(I_s+\delta V^*D^{-1}V)
=\det D\det\bigl(V^*(I+\delta D^{-1})V\bigr),
\]
and Cauchy--Binet proves \eqref{projexp}. Polynomial continuation removes invertibility. The same formula holds for $Q$ with probabilities $q_Y$.

Assume first $b_1\ne b_0$, and write
\[
\delta=b_1-b_0,\quad \tau=\frac{b_2-b_0}{b_1-b_0},\quad
UD_bU^*=b_0I+\delta P+\tau\delta vv^*.
\]
The determinant is affine in a rank-one update, so
\begin{equation}\label{affine}
\det(D_a+UD_bU^*)=(1-\tau)\det(D_a+b_0I+\delta P)
+\tau\det(D_a+b_0I+\delta Q).
\end{equation}
The scalars $1-\tau,\tau$ need not be nonnegative. They are not the proposed convex weights. Substitute \eqref{projexp} into \eqref{affine} using the single coupling $\mu$ from Lemma~\ref{couple}. If $Y\setminus X=\{j\}$, the paired summand is exactly
\begin{align*}
&(1-\tau)\prod_{i\in X}(a_i+b_1)\prod_{i\notin X}(a_i+b_0)
+\tau\prod_{i\in Y}(a_i+b_1)\prod_{i\notin Y}(a_i+b_0)\\
&\hspace{8pt}=\prod_{i\in X}(a_i+b_1)(a_j+b_2)\prod_{i\notin Y}(a_i+b_0).
\end{align*}
This is an actual permutation product for the original multiset $b$. Its coefficient is $\mu_{X,Y}\ge0$, and all such coefficients sum to one. Assign any bijections inside the repeated-value classes to obtain $\lambda_\pi$. The weights depend only on the two projection laws, hence only on $U,H,c$. Since the final identity is polynomial in all three values and has fixed weights, it also holds when values coincide. This proves Theorem~\ref{three}.

\subsection{A fully exact five-dimensional example}
Let $H(w)=I-2ww^T/(w^Tw)$ and set
\[
U=H(1,2,3,4,5)H(2,-1,1,0,3)
=\frac1{825}\begin{pmatrix}
499&88&-238&-120&-594\\448&451&74&-240&462\\
122&-286&661&-360&-132\\456&-528&-72&345&264\\
-90&-330&-420&-600&165
\end{pmatrix}.
\]
Take $H=\{1,2\}$, $c=3$, and $L=\{4,5\}$. A rational coupling is printed below. Entries are the integer numerators of $\mu_{X,Y}$ with common denominator $680625$; omitted admissible edges have weight zero.

\begin{center}\begin{tabular}{rrr@{\qquad}rrr}
\toprule $X$&$Y$&Numerator&$X$&$Y$&Numerator\\\midrule
12&123&50625&23&235&27200\\
13&123&596&24&124&26944\\
13&135&34000&24&234&202752\\
14&124&1280&24&245&57600\\
14&134&82944&25&235&16900\\
14&145&51200&34&345&6400\\
15&125&3600&35&135&6400\\
15&135&32500&45&145&6400\\
23&123&16900&45&345&51200\\
23&234&5184&&&\\
\bottomrule\end{tabular}\end{center}

The accompanying standard-library Python verifier checks $U^TU=I$, both marginals, nonnegativity, total mass one, all 32 coefficients in the variables $a_i$, and all 132 resulting scalar coefficients in $b_0,b_1,b_2$, using exact fractions. As a separate Gaussian-rational check, for
\[
a=(1+i,2-i,-1+2i,3+i,-2-i),\quad
b=(-2+i,-2+i,3+3i,1-2i,1-2i),
\]
both sides equal
\[
\frac{894536}{226875}-\frac{3793184}{75625}i.
\]

\section{A complete small-spectrum theorem in every dimension}
\begin{theorem}\label{small}
Fix $n\ge2$, nonzero complex numbers $a_i$, and complex numbers $\beta_j$. Put $x_i=a_i^{-1}$ and
\[
s_\pi=\sum_i x_i\beta_{\pi(i)}.
\]
Assume the $n!$ points $s_\pi$ are distinct and no three are collinear. Then there exists $\varepsilon>0$ such that for every $U\in U(n)$ and every complex $t$ with $|t|<\varepsilon$,
\[
\det(D_a+tUD_\beta U^*)\in\conv\left\{\prod_i(a_i+t\beta_{\pi(i)}):\pi\in S_n\right\}.
\]
The hypothesis holds on an open dense subset of $(\C\setminus\{0\})^n\times\C^n$.
\end{theorem}

\subsection{The limit and all its equality cases}
Put $d=\prod_i a_i$ and extend
\[
G_t(U)=\frac{d_{a,t\beta}(U)/d-1}{t}
\]
polynomially to $t=0$. The minor expansion gives
\begin{equation}\label{Gt}
G_t(U)=\sum_{k=1}^n t^{k-1}\sum_{|I|=|J|=k}x_I\beta_J|\det U[I,J]|^2,
\qquad G_0(U)=\sum_{ij}x_i\beta_j|U_{ij}|^2.
\end{equation}
For $t\ne0$ this complex affine normalization preserves convex-hull membership.

Fix a unit support phase $\eta$. Write $F_{\eta,t}=\Re(\eta G_t)$ and $h=\max_\pi\Re(\eta s_\pi)$. Since $|U|^2$ is doubly stochastic, $F_{\eta,0}\le h$ by the Birkhoff theorem. At most two permutations maximize, by the no-three-collinear assumption.

If two maximize, their relative permutation has exactly one nontrivial cycle. Indeed, splitting two or more disjoint cycles gives intermediate permutations whose nonnegative deficits sum to zero, producing at least three maximizers. Let $E$ be the union of the optimal permutation supports and
\[
L_E(U)=\sum_{(i,j)\notin E}|U_{ij}|^2.
\]
Let $\delta>0$ be the smallest nonoptimal permutation deficit. Any Birkhoff decomposition of $|U|^2$ shows
\begin{equation}\label{leak}
h-F_{\eta,0}(U)\ge\frac{\delta}{n}L_E(U).
\end{equation}
In fact a nonoptimal permutation uses at most $n$ forbidden edges, while each has deficit at least $\delta$. If no nonoptimal permutations exist, $n=2$ and the entire orbit is the already solved two-dimensional block.

The equality set $S$ is now explicit. One maximizer gives its monomial phase torus. Two differing by a transposition give a full $U(2)$ block and scalar phase blocks elsewhere. Two differing by a cycle of length at least three give only the two monomial tori: a strictly intermediate doubly stochastic matrix on that cycle would have two successive rows overlapping in exactly one nonzero column, impossible for a unitary.

\subsection{Quadratic coercivity and persistence of normal stationarity}
Use the bi-invariant Frobenius metric on $U(n)$. At $U_0\in S$, a skew-Hermitian generator $K$ gives
\[
L_E(e^{uK}U_0)=u^2\sum_{(i,j)\notin E}|(KU_0)_{ij}|^2+O(u^3).
\]
The quadratic kernel is precisely the tangent space of $S$. For a monomial torus it consists of diagonal generators. For the $U(2)$ block it consists of block-diagonal generators. For a longer cycle, any additional allowed off-diagonal entry has a forbidden transpose, so skew-Hermiticity again leaves only diagonal generators. Compactness of $S$ and its unit normal bundle therefore gives a uniform strictly negative normal Hessian for $F_{\eta,0}$.

At every point of $S$, $D_a+tU_0D_\beta U_0^*$ is block diagonal. A normal generator is off-block, so its first variation is off-block. The adjugate is block diagonal, even in singular cases by polynomial continuation. The determinant derivative is consequently zero in every normal direction, for every complex $t$ and every support phase. Thus the perturbation
\[
E_t=G_t-G_0
\]
has both a factor $t$ and zero first normal derivative on $S$.

\subsection{Uniform perturbation proof}
In a sufficiently small normal tube, write $U=e^KU_0$ with $K$ normal at $U_0$. Uniform Taylor estimates yield constants $c,C,C'>0$ such that
\[
F_{\eta,0}(e^KU_0)-F_{\eta,0}(U_0)\le-c\|K\|_F^2,
\]
\[
|E_t(e^KU_0)-E_t(U_0)|\le C|t|\|K\|_F^2.
\]
Changing the support phase from $\eta$ to $\eta'$ contributes at most $C'|\eta'-\eta|\|K\|_F^2$. These are quadratic estimates, not merely uniform continuity. For sufficiently small $|t|$ and $|\eta'-\eta|$,
\[
F_{\eta',t}(U)\le F_{\eta',t}(U_0).
\]
On a monomial torus the right side is a permutation value. On a $U(2)$ block the determinant is exactly a convex combination of its two permutation endpoints, with weights $|u_{11}|^2$ and $|u_{12}|^2$. Multiplication by the other scalar block determinants preserves that identity. Thus the desired bound holds throughout the tube.

Outside a smaller tube the compact unitary group has a positive base gap; uniform continuity preserves it for small parameters. This proves the support inequality for all $U$, for every $\eta'$ in a neighborhood of $\eta$, and all sufficiently small complex $t$. Finitely many phase neighborhoods cover the unit circle. Taking the minimum of their positive radii proves Theorem~\ref{small}. The original determinant at $t=0$ is just $d$.

\subsection{Genericity is part of the proof}
Every difference $s_\pi-s_\rho$ is a nonzero bilinear polynomial in $x,\beta$. If the collinearity polynomial for three distinct permutations vanished identically, the holomorphic ratio
\[
\frac{s_\rho-s_\pi}{s_\tau-s_\pi}
\]
would be real-valued on a complex open set where its denominator is nonzero, hence a real constant. Comparing coefficients would give
$P_\rho-P_\pi=c(P_\tau-P_\pi)$ for permutation matrices. Three distinct vertices of a polytope cannot be collinear. Hence every excluded equality or collinearity is a proper real-algebraic condition. Their finite union has empty interior, proving the open dense assertion. Uniformity in a sufficiently small parameter neighborhood follows from the same finite-cover argument.

\subsection{An exact genuinely complex instance}
Take
\begin{align*}
x&=(16+27i,\ 14-7i,\ 16+20i,\ -19-27i),\\
\beta&=(-13+8i,\ -22+11i,\ 6+25i,\ -22+6i),\qquad a_i=1/x_i.
\end{align*}
The 24 assignment sums are distinct. All $\binom{24}{3}=2024$ integer orientation determinants are nonzero; their minimum absolute value is $315$. The accompanying verifier checks these assertions with Gaussian integers, not floating-point geometry.
\subsection{Explicit radius: a finite quantitative certificate}
For the displayed example, one certified radius is
\[
\boxed{|t|\le10^{-22}.}
\]
This conservative value is proved by the following explicit estimates; it is not a numerical search threshold. Both spectra are nonconcyclic when $t\ne0$: the circle-test determinants with rows $(|z|^2,\Re z,\Im z,1)$ are $173208$ for $x$ and $36750$ for $\beta$. Inversion carries generalized circles to generalized circles, so $a=1/x$ is also nonconcyclic.

Here are sufficient computable conditions for $n\ge3$. Let $\alpha>0$ be the smallest absolute triple orientation among the $s_\pi$, let $D$ bound their diameter, and let $\ell>0$ be a lower bound on their pairwise distances. Choose upper bounds
\[
A_k\ge e_k(|x|)e_k(|\beta|),\quad B_k=4k^2A_k,\quad
S=\sum_{k=2}^nA_k,\quad T=\sum_{k=2}^nB_k,
\]
\[
Q=4B_1S/\ell+T,\qquad C\ge\frac{\pi^2}{2}(1+12n).
\]
Every $0<r\le1$ satisfying
\begin{equation}\label{radius}
4DSr+4S^2r^2<\alpha/2,\qquad CrQ<\frac{\alpha}{nD}
\end{equation}
is an admissible complex radius.

To prove this, first strengthen \eqref{leak} with the global distance estimate
\begin{equation}\label{globaldist}
\operatorname{dist}(U,S)^2\le C L_E(U).
\end{equation}
For a monomial or $U(2)$ block equality family, mask $U$ to its allowed blocks. Polar decomposition within each block gives $V\in S$ with
\[
\|U-V\|_F^2=L_E(U)+\sum_j(1-\sigma_j)^2
\le L_E(U)+\sum_j(1-\sigma_j^2)=2L_E(U),
\]
where the block singular values lie in $[0,1]$. The geodesic distance is at most $\pi/2$ times the chord distance.

For a cycle of length $k\ge3$, relabel its allowed entries to $(i,i)$ and $(i,i+1)$ cyclically, and put $d_i=|U_{ii}|^2$, $c_i=|U_{i,i+1}|^2$. Write $L=L_E(U)$. Row and column normalization imply $|c_i-c_j|\le2L$. Orthogonality of consecutive rows and Cauchy--Schwarz give $c_i d_{i+1}\le2L$, since all terms other than their one shared allowed column use a forbidden entry. Thus with $c=c_1$, $d_2\ge1-c-3L$, and
\[
c(1-c)\le5L,\qquad \min(c,1-c)\le10L.
\]
Consequently either all $c_i\le12L$, or all $d_i\le12L$. One of the two allowed permutations then has total off-permutation mass at most $(1+12n)L$. Choosing phases at its selected entries gives chord distance squared at most $2(1+12n)L$, proving \eqref{globaldist}.

Exterior powers give a uniform geodesic second derivative estimate
\[
\left|\frac{d^2}{du^2}|\det(e^{uK}U)[I,J]|^2\right|
\le4k^2\|K\|_F^2\quad(|I|=|J|=k).
\]
Indeed the induced exterior generator has norm at most $k\|K\|_F$; differentiate the squared matrix coefficient twice. By \eqref{Gt}, the Hessian norm of
$H=F_{\eta,t}-F_{\eta_0,0}$ is at most
\[
M=B_1|\eta-\eta_0|+\sum_{k=2}^nB_k|t|^{k-1}.
\]
Take a closest point $V\in S$ and a minimizing geodesic from it to $U$. Its initial direction is normal to $S$, so $H$ has zero first derivative there. Integration of the second derivative yields
\[
H(U)-H(V)\le\tfrac12 M\operatorname{dist}(U,S)^2\le\tfrac12 MC L_E(U).
\]
Together with \eqref{leak}, the support bound holds whenever $MC/2\le\delta/n$.

Finally, every normalized permutation value $G_t(P_\pi)$ moves by at most $R=rS$. The change of any triple orientation is at most $4DR+4R^2$. The first condition in \eqref{radius} therefore preserves all supporting edges and their labels. At an original edge, each nonendpoint gap is at least $\alpha/D$. Its perturbed outward unit normal differs by at most $4R/\ell$ from the original one. Thus $M\le rQ$, and the second condition in \eqref{radius} proves the support bound at every perturbed polygon edge. Their half-planes intersect in the desired convex hull.

For the exact example use $X=(43,21,36,46)$ and $Y=(21,33,31,28)$ as integer upper bounds on $|x|,|\beta|$. This gives
\[
(A_1,A_2,A_3,A_4)=(16498,37059829,15799252770,899499740832),
\]
\[
S=915336053431,\quad Q=241677564701544440,
\quad \alpha=315,\quad D=3764,\quad\ell=1,\quad C=245.
\]
The inequality $C>(\pi^2/2)49$ follows already from $\pi^2<10$. With $r=10^{-22}$, the two left sides in \eqref{radius} are respectively less than $1.379\cdot10^{-6}$ and $0.005922$; their right sides are $157.5$ and greater than $0.02092$. The verifier checks the unrounded rational comparisons exactly.

The same example admits the stronger certified radius $10^{-12}$ by removing its common degree-four translation and estimating the nine polygon edges separately. Appendix C of the bundled \nolinkurl{variational_notes.txt} supplies the complete bounds; \nolinkurl{variational_improved_radius_check.py} verifies them with rational arithmetic. The conservative calculation above is self-contained in this report.


\section{An explicit order-four cone and its sharp budget}
For equal-size subsets, write $\rho_{I,J}(U)=|\det U[I,J]|^2$. A common-minor cost is a real linear combination of these functions. It is determined by its values on permutations: the exact identity
\begin{equation}\label{omega}
g(U)=\sum_{\pi\in S_n}\Re\omega_\pi(U)g(P_\pi),\qquad
\omega_\pi(U)=\frac{\operatorname{sgn}(\pi)\prod_i U_{i,\pi(i)}}{\det U}
\end{equation}
follows by Laplace expansion and complementary minors. The weights sum to one but need not be nonnegative.

\begin{theorem}\label{cone}
For $n=4$, suppose a common-minor cost $g$ is nonnegative on all permutations and satisfies $g(\pi\circ(12))=g(\pi)$. Then $g$ is a nonnegative combination of singleton and two-minor indicators, hence $g(U)\ge0$ for all complex unitaries.
\end{theorem}
\begin{proof}
Write $g_{jk}=g(\pi)$ when $j=\pi(3)$ and $k=\pi(4)$. Averaging the minor indicators under $(12)$ shows that the invariant space consists precisely of arrays with
\[
g_{jk}-g_{kj}=d_j-d_k.
\]
Indeed symmetric edge indicators and row-three singleton indicators give sufficiency. For necessity, a mixed two-row indicator averages to one half of a directed cut $1_{j\in S,k\notin S}$; the other indicators are symmetric or singleton functions. Complementary minors handle larger sizes.

Set $m_{jk}=\min(g_{jk},g_{kj})\ge0$. Then
\[
g_{jk}=m_{jk}+(d_j-d_k)_+.
\]
Ordering the four potentials expresses the latter term as a nonnegative sum of threshold cuts. Thus there are six symmetric-edge generators $E_{ab}$ and fourteen nonempty proper directed-cut generators $C_S$. They have positive minor extensions
\[
E_{ab}=\rho_{34,ab},\qquad C_{\{a\}}=\rho_{3,a},\qquad
C_{[4]\setminus\{a\}}=\rho_{4,a},\qquad
C_{\{a,b\}}=\rho_{13,ab}+\rho_{23,ab}.
\]
These identities hold on permutations and therefore everywhere by \eqref{omega}.
\end{proof}

The theorem is a coefficient result with no spectral assumptions. Its full-support conclusion is also recoverable from N056/N059 of the supplied volume, and repeated-eigenvalue MOC itself is already known. It is recorded as a recovered theorem with an explicit cone proof; the following sharp budget and quantified neighborhood are the refinements developed here.

For the isoclinic family from N093/N129 of the supplied volume, put
\[
J=\begin{pmatrix}0&-1&-1&-1\\1&0&-1&1\\1&1&0&-1\\1&-1&1&0\end{pmatrix},
\quad U(r)=\frac{rI+J}{\sqrt{r^2+3}}.
\]
If $A=g(\id)$, $T$ is the sum of the six transposition $g$-values, $D$ the sum of the nine derangement $g$-values, and
$K=-\sum_{\pi\text{ a 3-cycle}}\bigl(\prod_{i:\pi(i)\ne i}J_{i,\pi(i)}\bigr)g(\pi)$, then
\[
g(U(r))=\frac{Ar^4+Tr^2-Kr+D}{(r^2+3)^2}.
\]
Every generator above has $\left(\begin{smallmatrix}2T&K\\K&2D\end{smallmatrix}\right)\succeq0$, as the exact table shows.
\begin{center}\begin{tabular}{lrrr@{\qquad}lrrr}
\toprule Generator&$T$&$D$&$K$&Generator&$T$&$D$&$K$\\\midrule
$E_{12}$&0&4&0&$C_{12}$&2&2&0\\
$E_{13}$&1&1&2&$C_{13}$&3&2&-2\\
$E_{14}$&1&1&-2&$C_{14}$&1&5&2\\
$E_{23}$&1&1&-2&$C_{23}$&3&2&2\\
$E_{24}$&1&1&2&$C_{24}$&1&5&-2\\
$E_{34}$&2&1&0&$C_{34}$&2&2&0\\
$C_1$&1&3&0&$C_{123}$&3&0&0\\
$C_2$&1&3&0&$C_{124}$&1&3&0\\
$C_3$&3&0&0&$C_{134}$&1&3&0\\
$C_4$&1&3&0&$C_{234}$&1&3&0\\\bottomrule
\end{tabular}\end{center}
Positive combinations preserve this matrix inequality. Hence $K^2\le4TD$ throughout this cone. Equality for $E_{13}$ proves sharpness of the constant. Relabeling and inversion give all row/column transpositions and corresponding isoclinic orientations; the sign of $K$ may reverse, without affecting the bound.

\subsection{A checked neighborhood without the symmetry}
The exact verifier in the research bundle also establishes the following finite strengthening. Let $g_0$ be as in Theorem~\ref{cone}, with zeros exactly at $\id,\tau=(12)$, and let
$m=\min_{\pi\notin\{\id,\tau\}}g_0(\pi)>0$. If $g$ is a common-minor cost with
\[
g(\id)=0,\quad g(\pi)\ge0,\quad \|g-g_0\|_\infty\le m/84,
\]
then $g$ still has a nonnegative-minor certificate.

Here is the exact finite justification. Among the 34 distinct singleton/two-minor indicators, 23 vanish at both $\id,\tau$. Their sum $H$ is $\tau$-invariant and $\max_\pi H(\pi)=6$. Thus $g_0-(m/6)H$ is in the cone of Theorem~\ref{cone}. Twelve of these indicators form a basis of the subspace vanishing at both points. Adding $\rho_{1,2}$ gives a basis of costs vanishing at $\id$. The rational inverse evaluation matrix has row $\ell^1$ norm at most 14 in its first twelve rows; the thirteenth coefficient is $g(\tau)\ge0$. Therefore any negative perturbation coefficient has magnitude at most $14\|g-g_0\|_\infty\le m/6$ and is absorbed by the corresponding reserve in $(m/6)H$. The script prints the selected minors, evaluation permutations, and the complete inverse matrix, and verifies them using exact fractions.

For a concrete spectral support consequence, take real $|u|\le1/200$ and
\[
a=(u+i,-u+i,1,2+i),\qquad b=(-1+i,-2-2i,2-i,1+2i).
\]
The checked gap at $u=0$ has height 66, zeros only at $\id,(12)$, and all other gaps at least 10. Direct expansion gives, for $u\ge0$, height $66+6u-12u^2$ and gap perturbation at most $22u+16u^2\le69/625<10/84$. Negative $u$ follows by swapping the first two rows. Consequently
\[
\Re\det(D_a+UD_bU^*)\le66+6|u|-12u^2\qquad(U\in U(4)).
\]
This is one complete support inequality, not all support directions for this spectral family.

\section{Full MOC for an explicit complex spectral family}\label{fullfamily}
\begin{theorem}
For every real $u$ with $|u|\le1/20$, let
\[
a(u)=(u+i,-u+i,1,2+i),\qquad
b=(-1+i,-2-2i,2-i,1+2i).
\]
Then $d_{a(u),b}(U)\in C(a(u),b)$ for every $U\in U(4)$. For $u\ne0$, both spectra are simple and noncollinear.
\end{theorem}

The proof is an exact polynomial certificate on the whole interval. All its coefficients are supplied in \nolinkurl{moc4_full_family_certificate.json}; the independent verification-only program is \nolinkurl{moc4_full_family_verify.py}. Both are included in the research bundle. No optimization solver or sampling enters the verification.

For $0<u\le1/20$, the permutation polygon has the following nine vertices in counterclockwise order, written as one-based permutation labels:
\[
3412,\quad4312,\quad3241,\quad4231,\quad1234,\quad
2134,\quad2143,\quad4123,\quad4132.
\]
Every vertex $z_\pi(u)$ is a complex polynomial of degree at most two. For consecutive labels $l,r$, including the closing pair, define
\[
G_{lr}(w;u)=\Im\bigl(\overline{z_r(u)-z_l(u)}(w-z_l(u))\bigr).
\]
The saved certificate gives, for every one of the nine edges, the identity
\begin{equation}\label{fullcert}
G_{lr}(d_{a(u),b}(U);u)=
\sum_{v=1}^{12}\lambda_{lr,v}(u)|\det U[I_{lr,v},J_{lr,v}]|^2,
\end{equation}
where each minor has size one or two and every $\lambda$ is a rational polynomial of degree at most four. The 108 coefficient polynomials are completely printed in the JSON, in increasing power order. The checker reconstructs the original spectral products and verifies every coefficient in \eqref{fullcert} on all 24 permutations. Identity \eqref{omega} then proves \eqref{fullcert} for every unitary.

For $E=1/20$ and $\lambda(u)=\sum_{k=0}^4a_ku^k$, its degree-four Bernstein coefficients on $[0,E]$ are
\[
B_j=\sum_{k=0}^j a_k E^k\frac{\binom jk}{\binom4k},\qquad j=0,\ldots,4.
\]
The checker proves all these rational numbers nonnegative, and hence
\[
\lambda(Et)=\sum_{j=0}^4 B_j\binom4j t^j(1-t)^{4-j}\ge0\quad(0\le t\le1).
\]
Thus all nine left-hand sides of \eqref{fullcert} are nonnegative on the entire parameter interval for every $U$.

To certify that these nine inequalities are sufficient, the checker also verifies all 36 squared distances between listed vertices and all nine consecutive edge cross products. Their Bernstein coefficients are nonnegative, at least one is positive, and their value at $E$ is positive. They are therefore strictly positive for $0<u\le E$. Each edge supports every permutation point by \eqref{fullcert}; distinct vertices and strict turns show these exposed edges describe exactly the finite permutation hull. Intersecting their half-planes proves the theorem for $u>0$.

At $u=0$, determinant continuity and compactness of the finite convex-weight simplex give the limit inclusion. Negative $u$ interchanges the first two $a$-labels, preserving the universal unitary assertion. This proves the full closed interval. The verifier and its interpretation were checked independently by a second agent and rerun by the primary agent.

\section{A sharper uniform bound in order four}
\begin{proposition}
For every $U\in U(4)$,
\[
\sum_\pi|\Re\omega_\pi(U)|\le\sum_\pi|\omega_\pi(U)|=\per|U|\le\frac32.
\]
\end{proposition}
\begin{proof}
For nonnegative unit vectors $x,y\in\R^4$, set $v_{ij}=x_i y_j+x_j y_i$ for $i<j$. Then
\[
\sum_{i<j}v_{ij}^2=1+(x\cdot y)^2-2\sum_i x_i^2y_i^2
\le1+\tfrac12(x\cdot y)^2\le\tfrac32.
\]
The first inequality uses $(\sum_i x_i y_i)^2\le4\sum_i x_i^2y_i^2$. Expand the permanent across the first two and last two rows of $|U|$, and apply Cauchy--Schwarz to their six pair coordinates, with complementary column pairs. Each pair-vector norm is at most $\sqrt{3/2}$.
\end{proof}

Since the real weights in \eqref{omega} sum to one, their negative mass is at most $1/4$. Thus in order four,
\[
d_{a,b}(U)\in\frac54 C(a,b)-\frac14 C(a,b).
\]
For a real common-minor cost with permutation range $[m,M]$ this also gives
\[
m-\frac{M-m}{4}\le g(U)\le M+\frac{M-m}{4}.
\]
The coefficient bound $3/2$ is attained by the raw signed expansion at a real normalized Hadamard matrix. This does not claim sharpness of the convex enclosure: alternative representations can reduce the signed mass. In N104, equations (27)--(30), every use of the bound $24$, or its proposed replacement $\sqrt{24}$, can be replaced by $3/2$. This improves that estimate but does not make its error zero.


\begin{center}\begin{tabular}{p{0.43\textwidth}p{0.45\textwidth}}
\toprule Verification & Result \\\midrule
Nested projection proof & Independently checked exterior Hall inequalities and direct rank-one determinant interpolation.\\
Five-dimensional certificate & 19 positive rational weights; 32 $a$-coefficients and 132 scalar coefficients verified exactly.\\
Small-spectrum proof & Two independent agent reviews checked the equality sets, normal Hessian, perturbation estimate and phase uniformity.\\
Generic example & 24 distinct Gaussian-integer sums; all 2024 triple orientations nonzero; minimum absolute value 315.\\
Order-four cone & Exact rank 9; all 20 budget matrices PSD; relabeling signs independently checked.\\
Nonsymmetric neighborhood & Exact 23-indicator reserve, maximum 6, inverse bound 14, hence radius $m/84$.\\
Full complex family & Nine exact facet identities, 108 interval-positive coefficient polynomials, 36 positive vertex distances and nine strict turns.\\\bottomrule
\end{tabular}\end{center}

\section{Sources and reproducibility}
The supplied \emph{MOC Complete Standalone Volume} (17 October 2026) was the primary prior-work record, particularly its purpose and notation (pp.~12--13), N013, N029--N030, N039, N059, N093, N104, N129, N139, and N150. The attached \texttt{SKILL(4).md} and \texttt{metaproof(9).txt} supplied methodological constraints, not additional mathematical axioms.

The research bundle contains this PDF, its LaTeX source, the exact verifiers and certificate data, and the principal proof/review notes. Uncertified numerical experiments are omitted. Run the exact scripts with Python 3; their core checks use the standard library. The result is reproducible without treating a solver's feasibility flag as proof.

The recent primary literature check found a July 2026 paper still explicitly treating unrestricted MOC as open \cite{mulero}. Shitov's 2025 theorem concerns a Hermitian summand and at most three nonreal eigenvalues of the other summand \cite{shitov}. These are context, not assumptions filling any gap in the new arguments.

\begin{thebibliography}{9}
\bibitem{lyons} R. Lyons, \emph{Determinantal probability measures}, Publications Math\'ematiques de l'IH\'ES 98 (2003), 167--212. The finite adjacent-rank coupling is related to Theorem 6.2. \url{https://rdlyons.pages.iu.edu/pdf/bases.pdf}.
\bibitem{shitov} Y. N. Shitov, \emph{A further sufficient condition for the determinantal conjecture}, Izvestiya: Mathematics 89(4) (2025), 862--869. \url{https://www.mathnet.ru/eng/im9292}.
\bibitem{mulero} J. I. Mulero-Mart\'inez, \emph{A variational framework for determinantal inequalities of normal matrices: Successes and obstructions}, Linear Algebra and its Applications 740 (2026), 19--38. \url{https://doi.org/10.1016/j.laa.2026.03.019}.
\bibitem{cll} E. A. Carlen, E. H. Lieb and M. Loss, \emph{An inequality of Hadamard type for permanents}, arXiv:math/0508096. The order-four case used here is proved directly. \url{https://arxiv.org/abs/math/0508096}.
\end{thebibliography}
\end{document}
